\section{Fixed core periods and explicit wheel stabilization}
\label{periods:section}

The network reductions leave scalar integrals with a fixed Kirchhoff
polynomial and varying numerator and pole order.  Two different statements
are needed at this stage.  For an arbitrary fixed core, the full period
space is finite dimensional, but a numerical basis need not be known.  For
a wheel, the space can be identified explicitly, uniformly in the integer
pole order.  We prove both statements with the boundary conditions retained.
The graph-period interpretation is consistent with the relative pairs of
Bloch, Esnault, and Kreimer~\cite{bloch2006} and Brown's treatment of
polynomial numerators~\cite[Section~6.7]{brown2015}; the sector presentation
below makes the required relative relations explicit.

\subsection{Normalized projective forms}
\label{periods:normalization}

Let $G=(V,E)$ be a fixed connected graph, let $N=|E|\ge2$ be its number
of labeled edges, and let $h_G=|E|-|V|+1$ be its loop number.  Multiple
edges remain distinct.  Its resistance Kirchhoff polynomial is
\begin{equation}
 U_G(t)=\sum_{T\text{ spanning tree of }G}\prod_{e\notin T}t_e.
 \label{periods:kirchhoff}
\end{equation}
The complement $E\setminus T$ of a spanning tree is called a cotree.
Every term in \eqref{periods:kirchhoff} has degree $h_G$ and coefficient
one.  The projective differential form associated with the chosen edge
order is
\begin{equation}
 \Omega_G=\sum_{j=1}^{N}(-1)^{j-1}t_j\,
 dt_1\wedge\cdots\wedge\widehat{dt_j}\wedge\cdots\wedge dt_N,
 \label{periods:omega}
\end{equation}
where the hat denotes omission.  We orient the positive projective
simplex so that the induced measure in the gauge $\sum_e t_e=1$ is
positive.  With this orientation our scalar convention is
\begin{equation}
 I_q(G;\alpha)=
 \int_{\mathbb{RP}^{N-1}_{>0}}
 \frac{\prod_{e\in E}t_e^{\alpha_e-1}}{U_G(t)^q}\,\Omega_G,
 \qquad
 q\in\mathbb Z_{>0},\quad \alpha_e\in\mathbb Z_{>0},\quad
 \sum_e\alpha_e=q h_G.
 \label{periods:scalar}
\end{equation}
For positive integer arguments $\Gamma(k)=(k-1)!$; there is no division
by $\prod_e\Gamma(\alpha_e)$.  The homogeneity in
\eqref{periods:scalar} is precisely the condition for descent under
common positive rescaling of the Schwinger parameters.

A complete numerator is a homogeneous polynomial $P\in\mathbb Q[t_e:e\in E]$,
including arbitrary signs, with
\begin{equation}
 \omega_{P,q}=\frac{P(t)}{U_G(t)^q}\Omega_G,
 \qquad \deg P=q h_G-N.
 \label{periods:polynomial-form}
\end{equation}
If the required degree is negative, this polynomial space is zero.  All
integrability statements in this section mean absolute convergence.  We
allow finite sums of \eqref{periods:polynomial-form} at different integer
pole orders, while keeping the graph and its polynomial fixed.  The
one-edge cases are immediate: a bridge has no positive-weight family in
\eqref{periods:scalar}, and a self-loop has the constant period one.

For each $q$, let $\mathcal P_q(G)$ denote the rational vector space
spanned by every convergent scalar period \eqref{periods:scalar} at that
order.  This is the full scalar space; a prescribed physical numerator
family may occupy a proper subspace.

\subsection{Exact ordered sectors and integer convergence}
\label{periods:hepp}

Fix a permutation $\sigma$ of the edges.  The corresponding ordered
sector has $0<t_{\sigma(1)}<\cdots<t_{\sigma(N)}$, and we set the largest
parameter equal to one.  Introduce successive ratios by
\begin{equation}
 t_{\sigma(i)}=\prod_{j=i}^{N-1}x_j\quad(1\le i<N),
 \qquad t_{\sigma(N)}=1,
 \qquad 0<x_j<1.
 \label{periods:ratios}
\end{equation}
For a prefix $S_j=\{\sigma(1),\ldots,\sigma(j)\}$, write $h(S_j)$ for
its cycle rank, computed on the original vertex set.  Equivalently,
$|S_j|-h(S_j)$ is the number of edges in a maximal forest of that prefix.

\begin{lemma}[Greedy cotree factorization]
\label{periods:greedy}
The substitution \eqref{periods:ratios} gives the exact polynomial identity
\begin{align}
 U_G(t(x))&=\left(\prod_{j=1}^{N-1}x_j^{h(S_j)}\right)U_\sigma(x),
 \label{periods:factorization}\\
 U_\sigma(x)&=
 \sum_{T\text{ spanning tree}}
 \prod_{j=1}^{N-1}x_j^{|(E\setminus T)\cap S_j|-h(S_j)}.
 \label{periods:sector-polynomial}
\end{align}
Every exponent in \eqref{periods:sector-polynomial} is nonnegative, and
the constant coefficient is exactly one.  Consequently, on the closed cube,
\begin{equation}
 1\le U_\sigma(x)\le \#\{T:T\text{ is a spanning tree of }G\}.
 \label{periods:positive-denominator}
\end{equation}
\end{lemma}

\begin{proof}
For every spanning tree $T$, its intersection with a prefix is a forest.
Therefore
\begin{equation}
 |T\cap S_j|\le |S_j|-h(S_j),\qquad
 |(E\setminus T)\cap S_j|\ge h(S_j).
 \label{periods:cotree-bound}
\end{equation}
Run the greedy algorithm in the order $\sigma$, accepting an edge when
it joins two distinct current components.  After each prefix, the accepted
edges form a maximal forest of that prefix.  The final accepted edges
form a spanning tree whose cotree attains equality in every inequality
\eqref{periods:cotree-bound} simultaneously.  This is the standard greedy
description of Hepp sectors~\cite[Lemma~2.8]{panzer2023}.

The exponent of $x_j$ in a cotree monomial after
\eqref{periods:ratios} is its number of edges in $S_j$, proving
\eqref{periods:factorization}.  The greedy cotree contributes the constant
term.  It is unique: successive differences of the complete prefix-count
sequence determine membership of the first $N-1$ edges, and the fixed
cotree cardinality determines the last edge.  Finally, every monomial of
\eqref{periods:sector-polynomial} lies between zero and one on the cube,
and its constant monomial equals one.  This proves
\eqref{periods:positive-denominator} without replacing the full polynomial
by a selected monomial.
\end{proof}

The absolute Jacobian of \eqref{periods:ratios} is
$\prod_{j=1}^{N-1}x_j^{j-1}$.  A scalar form therefore becomes
\begin{equation}
 \frac{\prod_{j=1}^{N-1}
 x_j^{\sum_{e\in S_j}\alpha_e-qh(S_j)-1}}
 {U_\sigma(x)^q}\,dx_1\cdots dx_{N-1}.
 \label{periods:scalar-pullback}
\end{equation}
Every edge ordering is retained; the equal-parameter faces have zero
interior measure but will remain part of the relative boundary.

\begin{proposition}[Integer convergence, including complete numerators]
\label{periods:integer-convergence}
The scalar form \eqref{periods:scalar} converges if and only if
\begin{equation}
 \sum_{e\in S}\alpha_e>q h(S)
 \qquad\text{for every }\varnothing\ne S\subsetneq E.
 \label{periods:power-counting}
\end{equation}
For a homogeneous rational-coefficient polynomial $P$, the complete form
\eqref{periods:polynomial-form} converges if and only if every surviving
monomial of $P$, after exact collection of identical terms, satisfies
\eqref{periods:power-counting}.  Equivalently, every ordered-sector
pullback has a polynomial numerator over $U_\sigma^q$ and is regular
on the entire closed cube.
\end{proposition}

\begin{proof}
By \eqref{periods:positive-denominator}, integrability of
\eqref{periods:scalar-pullback} is equivalent to each exponent being
greater than $-1$.  Since the exponents are integers, they must be
nonnegative.  Every nonempty proper subset appears as a prefix of some
ordering.  This gives \eqref{periods:power-counting}, in agreement with
the positive power-counting characterization in
\cite[Proposition~2.9]{panzer2023}.

For the signed numerator, perform the complete substitution, factorization,
and Jacobian multiplication first.  The result is a finite Laurent
polynomial divided by $U_\sigma^q$.  A Laurent polynomial is a finite
sum of integer powers, which at this intermediate stage may be negative.
If the collected numerator has least exponent $r<0$ in $x_i$, it can be
written
\begin{equation}
 x_i^r\bigl(L(x_1,\ldots,\widehat{x_i},\ldots,x_{N-1})+x_i R(x)\bigr),
 \label{periods:leading-boundary}
\end{equation}
where $L$ is nonzero.  Choose a compact positive patch in the face $x_i=0$,
away from every other coordinate face, on which $|L|$ has a positive lower
bound.  Such a patch exists because a nonzero polynomial cannot vanish on
an open real set.  On a sufficiently small product neighborhood of this
patch, the absolute numerator is bounded below by a positive multiple of
$x_i^r$, while \eqref{periods:positive-denominator} bounds the denominator
above.  The absolute integral diverges because the integer $r$ is at most
$-1$.  Thus convergence removes every negative coordinate exponent.

Distinct homogeneous monomials cannot disappear through the substitution:
their prefix exponents determine the first $N-1$ original powers by
successive differences, and their common total degree determines the last
power.  The common Jacobian and denominator shifts preserve this
injectivity.  Hence every surviving original monomial has nonnegative
sector exponents.  Conversely, a finite sum of convergent scalar
monomials is absolutely convergent.  No cancellation between divergent
absolute integrals has been used.
\end{proof}

\subsection{A fixed relative boundary complex}
\label{periods:relative}

For one fixed sector, put $m=N-1$.  Writing $\mathbb A^m_{\mathbb Q}$
for affine $m$-space over the rational numbers, define
\begin{equation}
 A_\sigma=\mathbb Q[x_1,\ldots,x_m,U_\sigma^{-1}],\qquad
 X_\sigma=\mathbb A^m_{\mathbb Q}\setminus\{U_\sigma=0\}.
 \label{periods:ring}
\end{equation}
Thus $A_\sigma$ consists of rational functions with denominators that are
powers of the one fixed polynomial $U_\sigma$.  Let $B_\sigma$ be the
union, inside $X_\sigma$, of the $2m$ faces $x_i=0$ and $x_i=1$.  These
faces have simple normal crossings: locally, all components through a
point are distinct coordinate hyperplanes.  Every nonempty intersection
of compatible faces is smooth and affine.  The cube is a relative chain
in $(X_\sigma,B_\sigma)$ because of
\eqref{periods:positive-denominator}.

We describe its relative differential complex explicitly.  In degree $p$,
take the regular algebraic $p$-forms on $X_\sigma$ whose restriction to
every component of $B_\sigma$ vanishes.  If a coordinate wedge contains
precisely the differentials indexed by $J\subseteq\{1,\ldots,m\}$, its
coefficient must be divisible in $A_\sigma$ by
\begin{equation}
 \prod_{i\notin J}x_i(1-x_i).
 \label{periods:boundary-factors}
\end{equation}
For $i\in J$, the wedge already restricts to zero on both faces normal
to $x_i$, so no corresponding coefficient factor is required.

\begin{lemma}[Relative presentation]
\label{periods:relative-presentation}
The top relative algebraic de Rham group is the rational vector-space
quotient
\begin{equation}
 \boxed{
 H_{\mathrm{dR}}^m(X_\sigma,B_\sigma;\mathbb Q)
 =\frac{A_\sigma}
 {\displaystyle\sum_{i=1}^m
 \partial_{x_i}\bigl(x_i(1-x_i)A_\sigma\bigr)}.}
 \label{periods:relative-quotient}
\end{equation}
It is finite dimensional over $\mathbb Q$.
\end{lemma}

\begin{proof}
For each form degree, augment restriction to the individual faces by
alternating restriction to their double intersections, triple intersections,
and so on.  In local normal-crossing coordinates this finite face complex
is exact after adjoining the kernel
\eqref{periods:boundary-factors}.  The corresponding alternating complex
of constant sheaves resolves the constant sheaf on the union of the
faces: at a point it is the augmented simplex complex of the boundary
components through that point.  The algebraic de Rham comparison theorem
on $X_\sigma$ and on each smooth affine face
intersection~\cite[Theorem~1]{grothendieck1966} consequently identifies
the kernel differential complex with relative cohomology.

In customary notation the kernel is
$\Omega_{X_\sigma}^{\bullet}(\log B_\sigma)(-B_\sigma)$, the logarithmic
differential complex twisted to vanish on the boundary.  Formula
\eqref{periods:boundary-factors} describes its terms without relying on
that notation.  Its terms are locally free on the affine variety
$X_\sigma$, so their higher coherent cohomology vanishes and global
sections compute its cohomology.  In degree $m$ the global term is
$A_\sigma\,dx_1\wedge\cdots\wedge dx_m$.  In degree $m-1$, the term
omitting $dx_i$ has coefficient in $x_i(1-x_i)A_\sigma$.  Its differential
gives exactly the denominator of \eqref{periods:relative-quotient}.

The complement and every face intersection are algebraic varieties of
finite type, with finite-dimensional singular cohomology.  The finite
relative face complex and comparison therefore give finite-dimensional
de Rham cohomology after extension of scalars to $\mathbb C$.  The
complex is defined over $\mathbb Q$, and field extension commutes with
its cohomology; hence its dimension over $\mathbb Q$ is finite as well.
\end{proof}

The denominator of \eqref{periods:relative-quotient} is a vector
subspace, generally not an ideal.  Multiplication of a derivative relation
by a polynomial requires its product-rule correction.  For every
polynomial $P$, coordinate $i$, and integer $q\ge0$, the complete
relation is
\begin{align}
 \partial_{x_i}\left(\frac{x_i(1-x_i)P}{U_\sigma^q}\right)
 ={}&\frac{(1-2x_i)P+x_i(1-x_i)\partial_{x_i}P}{U_\sigma^q}
 \notag\\[-2pt]
 &-q\frac{x_i(1-x_i)P\,\partial_{x_i}U_\sigma}{U_\sigma^{q+1}}.
 \label{periods:complete-relation}
\end{align}
Every term is regular on the closed cube.  The primitive vanishes on the
two faces normal to $x_i$, and its differential wedge restricts to zero
on the other faces.  Its integral is therefore exactly zero.  Allowing
all polynomials and all pole orders in
\eqref{periods:complete-relation} gives the full relative-exact subspace
in \eqref{periods:relative-quotient}.  This completeness statement does
not identify a particular finite row bank with that subspace, and does
not assert that every numerical identity between periods is cohomological.

\subsection{Finite span and stabilization of the full scalar levels}
\label{periods:stabilization}

\begin{theorem}[Fixed core finite span and eventual stabilization]
\label{periods:finite-span}
For a fixed connected graph $G$, the rational span of every absolutely
convergent form \eqref{periods:polynomial-form}, at all positive integer
pole orders, is finite dimensional.  Moreover,
\begin{equation}
 \mathcal P_q(G)\subseteq\mathcal P_{q+1}(G),
 \qquad
 \mathcal P_q(G)=\bigcup_{r\ge1}\mathcal P_r(G)
 \quad\text{for all sufficiently large }q.
 \label{periods:stabilization-equation}
\end{equation}
The eventual threshold is existential; the assertion includes no
numerical basis or arithmetic-independence claim.
\end{theorem}

\begin{proof}
Proposition~\ref{periods:integer-convergence} puts every input into the
finite direct sum, over all edge orderings, of the fixed spaces
\eqref{periods:relative-quotient}.  Integration over each cube is a
rational linear functional on its relative cohomology.  Choose a rational
basis of each finite-dimensional group and regular top-form representatives.
Their cube integrals form a finite spanning list for all the original
periods.  Integration need not be injective, so this list need not be
numerically independent.  The artificial internal sector faces may
also make the direct sum larger than the original graph-family image.

For nesting, multiply the numerator by the complete Kirchhoff polynomial:
\begin{equation}
 I_q(G;\alpha)=
 \sum_{T\text{ spanning tree}}
 I_{q+1}(G;\alpha+\mathbf1_{E\setminus T}),
 \label{periods:dimension-shift}
\end{equation}
where $\mathbf1_{E\setminus T}$ adds one on exactly the cotree edges.
Every summand converges separately.  Its positive integrand is the old
positive integrand multiplied by
\begin{equation}
 0\le \frac{\prod_{e\notin T}t_e}{U_G(t)}\le1.
 \label{periods:dominated-shift}
\end{equation}
Equivalently, \eqref{periods:cotree-bound} preserves every strict
inequality in \eqref{periods:power-counting}.  This proves nesting.
Finally, choose a finite basis of the union from actual scalar periods.
Each of its finitely many elements occurs at a finite pole order.  The
maximum of these orders is a sufficient threshold by nesting.
\end{proof}

\begin{example}[A plateau does not locate the threshold]
\label{periods:plateau}
For the cycle with $N$ edges, including the two-edge cycle of parallel
edges, one has $h_G=1$ and $U_G=\sum_e t_e$.  There are no positive
integer weights with sum $q<N$.  In the unit-sum gauge, however,
\begin{equation}
 I_q(G;\alpha)=\frac{\prod_e(\alpha_e-1)!}{(q-1)!},
 \qquad
 I_N(G;1,\ldots,1)=\frac1{(N-1)!}.
 \label{periods:cycle-moments}
\end{equation}
Thus the full scalar spaces are zero at every $q<N$ and equal to
$\mathbb Q$ at every $q\ge N$.  Arbitrarily many consecutive equal
levels can precede a strict increase.  Finite dimension bounds the
number of strict increases, not the location of the final increase.
\end{example}

There is a constructive route to the finite relative space, though its
cost is a separate issue.  Algorithms for localization, restriction, and
de Rham cohomology of algebraic differential modules and their complexes
apply to the finite face complex~\cite{oaku2001,walther2000}.  Once
representatives $f_1,\ldots,f_r\in A_\sigma$ of a rational basis have
been computed, a requested integrand $f\in A_\sigma$ admits an identity
\begin{equation}
 f-\sum_{j=1}^r c_j f_j
 =\sum_{i=1}^m\partial_{x_i}\bigl[x_i(1-x_i)g_i\bigr],
 \qquad c_j\in\mathbb Q,\quad g_i\in A_\sigma.
 \label{periods:target-certificate}
\end{equation}
Searching increasing polynomial-degree and pole-order bounds for the
$g_i$ is finite rational linear algebra at each step, and existence of
the spanning basis guarantees eventual success.  The full rational
identity and its boundary factors are an independently checkable
certificate.  A bounded search against caller-supplied master forms is
useful without having computed a complete basis; its failure does not
certify independence or a nonzero residual value.

For example, when $U_\sigma=1+x$, partial fractions show that the relative
quotient is spanned by $1$ and $1/(1+x)$.  To see this with its boundary
condition, write a rational function as $f=a/(1+x)+R'$ with
$R\in\mathbb Q[x,(1+x)^{-1}]$.  Then
\begin{equation}
 f=R(1)-R(0)+\frac{a}{1+x}
 +\partial_x\bigl[R(x)-R(0)-x(R(1)-R(0))\bigr].
 \label{periods:one-dimensional-relative}
\end{equation}
The last primitive is divisible by $x(1-x)$ in that ring.  A rational
derivative has zero residue at $x=-1$; after removing that coefficient,
the endpoint condition also rules out a nonzero exact constant.  These
two classes are therefore independent in relative cohomology and have
periods $1,\log2$.  The original two-edge cycle nevertheless has only the
one-dimensional rational scalar span in \eqref{periods:cycle-moments}.
The sector cohomology can strictly exceed the original graph-family image.

The finite-span theorem applies coefficientwise to the complete finite
source polynomials, with every original free spinor index retained.
Changing an integer holomorphic dimension or taking finitely many current
insertions is covered whenever the preceding network reduction yields
the absolutely convergent normalized forms
\eqref{periods:polynomial-form} with the same $U_G$.  This does not make
all arbitrary complex source evaluations a finite-dimensional rational
space.  Dimension-dependent normalization factors must first be
separated.  Noninteger exponents, regulator logarithms, changing divisors,
and an unbounded family of different cores are outside the theorem.

\subsection{Arbitrary integer powers on a wheel}
\label{periods:wheel-upper}

Let $W_n$, with $n\ge3$, have hub $0$, cyclic rim vertices
$1,\ldots,n$, and spokes $0i$.  It has $2n$ edges and loop number $n$.
Write the spoke weights as $a_i$ and the rim weights on $i(i+1)$ as
$\rho_i$, with cyclic indices.  Then
$\sum_i(a_i+\rho_i)=nq$.  Convergence implies
$1\le a_i,\rho_i<q$ by applying
\eqref{periods:power-counting} to the complement of each edge, so $q\ge2$.

\begin{proposition}[Wheel upper inclusion]
\label{periods:wheel-upper-proposition}
Every convergent positive-integer wheel monomial belongs to
\begin{equation}
 I_q(W_n;\alpha)\in
 \mathbb Q+\sum_{j=2}^{n-1}\mathbb Q\,\zeta(2j-1).
 \label{periods:wheel-upper-space}
\end{equation}
Here $\zeta(k)=\sum_{s=1}^{\infty}s^{-k}$ for integers $k>1$.
\end{proposition}

\begin{proof}
We give the complete rational spectral reduction, including its
normalization.  Put $b_i=q-a_i$ and $\beta_i=q-\rho_i$ for the
corresponding position powers.  With $\xi_i\in\mathbb R^{2q}$, use the
Dirac delta to fix the first squared radius to one.  The projective
position integral is
\begin{equation}
 \left(\prod_i\Gamma(b_i)\Gamma(\beta_i)\right)
 \int\prod_i\frac{d^{2q}\xi_i}{\pi^q}\,
 \delta(|\xi_1|^2-1)
 \prod_i|\xi_i|^{-2b_i}|\xi_i-\xi_{i+1}|^{-2\beta_i}.
 \label{periods:position-wheel}
\end{equation}
It equals \eqref{periods:scalar}.  Indeed, Schwinger parametrization and
Gaussian integration give the inverse $q$th power of the grounded
conductance Laplacian determinant.  For conductances $s_e$ this determinant
is $(\prod_e s_e)U_G(s^{-1})$; inversion $s_e=t_e^{-1}$ changes each
power $\beta_e$ to $q-\beta_e$.  Removing the common logarithmic scale
by fixing one squared radius gives the delta factor in
\eqref{periods:position-wheel}.  In logarithmic radii $v_i=\log|\xi_i|$,
unit angular measure contributes $(2/\Gamma(q))^n$, and the fixed squared
radius contributes $1/2$.  Thus the radial prefactor before angular
eigenvalues is $2^{n-1}/\Gamma(q)^n$, multiplied by the displayed gamma
factors.

For angular harmonic degree $\nu\ge0$, set $s=\nu+q-1$.  The harmonic
space consists of restrictions to the unit sphere of homogeneous
harmonic polynomials of degree $\nu$.  Its dimension, obtained by
subtracting the traced polynomials, is
\begin{equation}
 M_q(s)=
 \binom{\nu+2q-1}{2q-1}-\binom{\nu+2q-3}{2q-1}
 =\frac{2s^2}{(2q-2)!}\prod_{j=1}^{q-2}(s^2-j^2).
 \label{periods:spherical-multiplicity}
\end{equation}
For a rim power $\beta$, put $d=q-1-\beta\ge0$.  Factoring out the
symmetric powers of the two radii, the angular eigenvalue at their ratio
$0<r<1$ is
\begin{equation}
 \sum_{k=0}^{d} E_{d,k}(s)r^{s+2k-d},\qquad
 E_{d,k}(s)=\frac{\Gamma(q)}{\Gamma(\beta)}
 \frac{(-1)^k\binom dk}{\prod_{j=0}^{d}(s+k-d+j)}.
 \label{periods:angular-eigenvalue}
\end{equation}
For completeness, if $(u)_k=u(u+1)\cdots(u+k-1)$ denotes a rising
factorial, angular integration of the binomial expansion gives
\begin{equation}
 r^{\nu+\beta}\frac{(\beta)_\nu}{(q)_\nu}
 \sum_{k=0}^{d}
 \frac{(\nu+\beta)_k(-d)_k}{(\nu+q)_k k!}\,r^{2k}.
 \label{periods:terminating-angular}
\end{equation}
Replacing the rising factorials by gamma ratios gives
\eqref{periods:angular-eigenvalue} term by term.  The unsplit angular
operator has nonnegative eigenvalues: the positive Laplace representation
of $(1-2rz+r^2)^{-\beta}$ is a positive mixture of $\exp(uz)$, $u>0$,
whose expansion consists of positive-semidefinite tensor kernels.
Consequently the cyclic harmonic traces are nonnegative and can be summed
under the convergent radial integral.

The radial drift at vertex $i$ is
$2q-2b_i-\beta_{i-1}-\beta_i$.  Its sum is zero; choose integers $c_i$
such that
\begin{equation}
 c_i-c_{i-1}=2q-2b_i-\beta_{i-1}-\beta_i,
 \qquad c_0=c_n=0.
 \label{periods:radial-drift}
\end{equation}
Selecting the term $k_i$ in each angular eigenvalue gives the radial
integrand
\begin{equation}
 \exp\left\{\sum_i(c_i-c_{i-1})v_i
 -\sum_i(s+2k_i-d_i)|v_i-v_{i+1}|\right\},
 \qquad v_1=0,\quad d_i=\rho_i-1.
 \label{periods:radial-integrand}
\end{equation}
In an ordering of the $v_i$, each gap has inverse decay coefficient
$\sum_{i\in S}(c_i-c_{i-1})+
\sum_{i\in\partial S}(s+2k_i-d_i)$, where $S$ is the set below the
gap and $\partial S$ denotes its crossing rim edges.  At $s=q-1$ and
all $k_i=0$, positivity follows by restricting the convergent position
integral to separated radial annuli and a fixed open angular set.
Increasing $s$ or any $k_i$ preserves it.  Hence each selected radial
term has a convergent gap integral.

The bilateral Laplace transform of $\exp(-p|v|)$ is
$2p/(p^2-w^2)$.  Summing a complete link before taking the convolution
residue uses the finite rational identity
\begin{equation}
 \sum_{k=0}^{d}
 \frac{(-1)^k\binom dk}{\prod_{j=0}^{d}(s+k-d+j)}
 \frac{2(s+2k-d)}{(s+2k-d)^2-w^2}
 =\frac{2\,4^d d!}
 {\prod_{k=0}^{d}((s+2k-d)^2-w^2)}.
 \label{periods:link-identity}
\end{equation}
Both sides are even rational functions of $w$ vanishing at infinity.
At a pole $w=-(s+2k-d)$ their residues agree because the product of the
other denominator differences is
\begin{equation}
 \prod_{j\ne k}\bigl((s+2j-d)^2-(s+2k-d)^2\bigr)
 =\frac{4^d(-1)^k k!(d-k)!}{s+2k-d}
 \prod_{j=0}^{d}(s+k-d+j).
 \label{periods:link-residues}
\end{equation}
The positive-pole residues then agree as well, proving
\eqref{periods:link-identity}.

For each link $i$, form the integers
$c_i-d_i,c_i-d_i+2,\ldots,c_i+d_i$.  They have one common parity by
\eqref{periods:radial-drift}.  Subtract the common minimum and divide
by two; let $N_j$ be the number of resulting occurrences of $j$, and let
$L$ be the largest occupied integer.  Gap positivity for every contiguous
rim segment and its complement implies
$|c_i-c_j|<\beta_i+\beta_j$.  Therefore
\begin{equation}
 \max_i(c_i+d_i)-\min_i(c_i-d_i)<2(q-1),
 \qquad 0\le L\le q-2.
 \label{periods:residue-strip}
\end{equation}
Each link contributes at most one occurrence to each $j$, so $N_j\le n$.

Define the univariate rational function and its principal coefficients by
\begin{align}
 P_*(w)&=\prod_{j=0}^{L}(w-j)^{-N_j}
       =\sum_{j=0}^{L}\sum_{k=1}^{N_j}\frac{A_{j,k}}{(w-j)^k},
 \label{periods:univariate-function}\\
 A_{j,k}&=\left.
 \frac{1}{(N_j-k)!}\left(\frac{d}{dw}\right)^{N_j-k}
 \prod_{l\ne j}(w-l)^{-N_l}\right|_{w=j}.
 \label{periods:principal-coefficients}
\end{align}
Zero multiplicities contribute no terms.  Let $N_{\mathrm{tot}}=\sum_jN_j$.
After \eqref{periods:link-identity}, rescaling the convolution variable
by two cancels the factor $2^{2N_{\mathrm{tot}}-1}$ from the radial
measure and links: the differential contributes $2$, and their
$N_{\mathrm{tot}}$ quadratic factors contribute $4^{N_{\mathrm{tot}}}$.
The remaining spectral rational function is
\begin{equation}
 F(s)=\sum_{j,k,l,t}
 \frac{(-1)^{N_{\mathrm{tot}}+t}
 \binom{k+t-2}{k-1}A_{j,k}A_{l,t}}
 {(s+l-j)^{k+t-1}},
 \label{periods:convolution-residues}
\end{equation}
where $1\le k\le N_j$ and $1\le t\le N_l$.  Indeed, this is the sum
of the left residues of $(-1)^{N_{\mathrm{tot}}}P_*(w)P_*(w-s)$;
expanding $(s+l-w)^{-t}$ at $w=j$ gives the displayed binomial coefficient.
All normalization factors now give the complete period
\begin{equation}
 I_q(W_n;\alpha)=
 \left(\prod_i(b_i-1)!d_i!\right)
 \sum_{s=q-1}^{\infty}M_q(s)F(s).
 \label{periods:wheel-residue-sum}
\end{equation}

Only integer poles $s=j-l$ occur, with $|j-l|\le q-2$.
Since $k,t\le n$, their order before multiplication by $M_q$ is at
most $2n-1$.  The multiplicity has a double zero at zero and a simple
zero at every other possible pole.  Thus the remaining order is at
most $2n-3$ at zero and at most $2n-2$ elsewhere.

The function $F$ is odd.  Translating $w$ exchanges the left residues
at $s$ with the right residues at $-s$; their sums are negatives because
$P_*(w)P_*(w-s)$ has no residue at infinity.  Since $M_q$ is even,
$M_q(s)F(s)$ is odd as well.  Opposite nonzero poles consequently have
coefficients $B_{-h,u}=(-1)^{u+1}B_{h,u}$ at order $u$.  Summing their
partial fractions over integers $s\ge q-1$ cancels every even zeta
coefficient.  At zero only odd pole orders occur.  There is no polynomial
part: the nonnegative harmonic series in
\eqref{periods:wheel-residue-sum} converges.  Its total simple-pole
coefficient also vanishes by convergence.  The finite lower-limit
corrections are rational harmonic sums.  The remaining zeta weights
are therefore exactly among $3,5,\ldots,2n-3$, proving
\eqref{periods:wheel-upper-space}.
\end{proof}

\subsection{Radial periods and the odd-zeta directions}
\label{periods:radial}

Set every spoke weight to $q'-1$ and every rim weight to one.  The
angular defects and radial drifts vanish.  The convolution of $n$
copies of $\exp(-s|v|)$ at zero is
\begin{equation}
 \int_{\mathbb R}\left(\frac{2s}{s^2+y^2}\right)^n\frac{dy}{2\pi}
 =2^{1-n}\binom{2n-2}{n-1}s^{1-n}.
 \label{periods:radial-convolution}
\end{equation}
Together with the radial measure and angular eigenvalues, this gives
\begin{equation}
 I_{q'}(W_n;q'-1,\ldots,q'-1;1,\ldots,1)
 =\binom{2n-2}{n-1}
 \sum_{s=1}^{\infty}
 \frac{2s^2\prod_{j=1}^{q'-2}(s^2-j^2)}
 {(2q'-2)!\,s^{2n-1}},
 \quad 2\le q'<n.
 \label{periods:radial-master}
\end{equation}
Terms below $s=q'-1$ vanish.  The tail converges precisely in the
displayed range.  At $q'=2$ this is
$\binom{2n-2}{n-1}\zeta(2n-3)$, the usual wheel period
\cite{bloch2006}.  As $q'$ increases from $2$ to $n-1$, the lowest zeta
weight in \eqref{periods:radial-master} is
$2n-2q'+1$, with coefficient
\begin{equation}
 \frac{2\binom{2n-2}{n-1}}{(2q'-2)!}\ne0.
 \label{periods:radial-diagonal}
\end{equation}
These values therefore span all the odd-zeta directions by a triangular
rational coefficient matrix.  This is a statement about explicit linear
combinations, not arithmetic independence of their numerical values.

\subsection{Exact star--triangle and parallel normalizations}
\label{periods:star-parallel}

We next construct a nonzero rational direction at the same pole level.
The normalization of the graph operations matters.  At a trivalent
vertex, let the incident resistance powers be $a_1,a_2,a_3$, with
$a_1+a_2+a_3=q$.  If $x_i$ is the resistance on the new triangle edge
opposite the old edge $i$, the positive transformation is
\begin{equation}
 r_i=\frac{x_jx_k}{x_1+x_2+x_3},\qquad
 U_{\triangle}=(x_1+x_2+x_3)U_{\star},\qquad
 \left|\frac{\partial(r_1,r_2,r_3)}{\partial(x_1,x_2,x_3)}\right|
 =\frac{x_1x_2x_3}{(x_1+x_2+x_3)^3}.
 \label{periods:star-map}
\end{equation}
The inverse is $x_i=(r_1r_2+r_2r_3+r_3r_1)/r_i$.  Substitution gives
$\prod_i x_i^{q-a_i-1}/U_\triangle^q$; the remaining power of
$x_1+x_2+x_3$ is $q-\sum_i a_i=0$.  Thus the opposite triangle powers
are $q-a_i$ and there is no scalar prefactor.  Homogeneity makes the
same identity valid for the projective measure.

For parallel resistances with powers $a,b$, use
$x=r/(1-z)$ and $y=r/z$, with $0<z<1$.  Since
$U_G=(x+y)U_{G'}(xy/(x+y))$ and
$|\partial(x,y)/\partial(r,z)|=r/[z^2(1-z)^2]$, the fiber integral is
\begin{equation}
 \int_0^1 z^{q-b-1}(1-z)^{q-a-1}\,dz
 =\frac{\Gamma(q-a)\Gamma(q-b)}{\Gamma(2q-a-b)},
 \qquad \alpha_{\mathrm{new}}=a+b-q.
 \label{periods:parallel-beta}
\end{equation}
All applications below satisfy $0<a,b<q$ and $a+b>q$.  The fiber is
therefore an ordinary convergent beta integral, and the new edge has a
positive integer power.  Positive changes of variables and these finite
fibers permit convergence to be checked on the final reduced graph.

\subsection{A nonzero rational coefficient for every pole order}
\label{periods:k4-direction}

On $K_4=W_3$, take spoke powers $(q-1,1,1)$ and cyclic rim powers
$(1,q-1,q-1)$, with $q\ge3$.  The three heavy edges form a path, as do
the complementary light edges.  Every triangle has power sum $q+1$ or
$2q-1$; every four-cycle has power sum greater than $q$; and removing
one edge leaves sum greater than $2q$.  These cover all potentially
nontrivial inequalities \eqref{periods:power-counting}, so the period
converges.

\begin{lemma}[Uniform rational direction on $K_4$]
\label{periods:k4-rational}
There is a rational number $R_q$ such that
\begin{equation}
 I_q(K_4;q-1,1,1;1,q-1,q-1)
 =R_q+\frac{12((q-2)!)^2}{(2q-2)!}\zeta(3),
 \qquad R_q\ne0\quad(q\ge3).
 \label{periods:k4-value}
\end{equation}
In fact $R_3=-1/6$, and $R_q>0$ for every $q\ge4$.
\end{lemma}

\begin{proof}
Put $r=q-2$.  The occupied multiplicities in
\eqref{periods:univariate-function} are $(2,\ldots,2,3)$ at
$0,\ldots,r$, and the factorial prefactor in
\eqref{periods:wheel-residue-sum} is $(r!)^4$.  Thus
\begin{equation}
 \begin{aligned}
 P_*(w)&=\frac1{(w-r)^3\prod_{j=0}^{r-1}(w-j)^2},\\
 A_{r,3}&=\frac1{(r!)^2},\\
 A_{j,2}&=-\frac1{(j!)^2((r-j)!)^2(r-j)}\quad(j<r).
 \end{aligned}
 \label{periods:k4-principal}
\end{equation}
Only the order-five pole at shift zero and the order-four poles at
nonzero shifts can contribute to $\zeta(3)$ after multiplication by
$M_q$.  Including $(r!)^4$, the first contribution is
$12(-1)^r(r!)^2/(2r+2)!$.  For paired shifts $k,-k$, $1\le k\le r$,
the relevant derivative is
\begin{equation}
 M_q'(k)=\frac{2(-1)^{r-k}k(r-k)!(r+k)!}{(2r+2)!}.
 \label{periods:multiplicity-derivative}
\end{equation}
Together with \eqref{periods:k4-principal} and the binomial factor three
in \eqref{periods:convolution-residues}, the paired contribution is
\begin{equation}
 \frac{12(r!)^2}{(2r+2)!}
 (-1)^{r-k}\binom rk\binom{r+k}{k}.
 \label{periods:paired-k4-contribution}
\end{equation}
The sum, including $k=0$, equals the coefficient in
\eqref{periods:k4-value}.  If $[x^r]$ denotes extraction of the
coefficient of $x^r$, the required identity is
\begin{equation}
 \sum_{k=0}^r(-1)^{r-k}\binom rk\binom{r+k}{k}
 =[x^r](1+x)^r((1+x)-1)^r=1.
 \label{periods:binomial-one}
\end{equation}
The upper inclusion already proved for $n=3$ leaves only a rational
coefficient $R_q$.

We specify the finite calculation needed below.  Define
$H_a^{(u)}=\sum_{k=1}^{a}k^{-u}$, with $H_0^{(u)}=0$.
The remaining principal coefficients are explicitly
\begin{align}
 A_{r,2}&=-\frac{2H_r^{(1)}}{(r!)^2},&
 A_{r,1}&=\frac{2(H_r^{(1)})^2+H_r^{(2)}}{(r!)^2},\notag\\
 A_{j,1}&=A_{j,2}\left(-\frac3{j-r}
 -2\sum_{\substack{0\le l<r\\l\ne j}}\frac1{j-l}\right)
 &&(0\le j<r).
 \label{periods:k4-lower-principal}
\end{align}
These follow by logarithmic differentiation of
\eqref{periods:principal-coefficients}.  Form $F(s)$ by
\eqref{periods:convolution-residues} using these complete coefficients.
Its product with
$M_q(s)$ has poles of order at most three:
\begin{equation}
 M_q(s)F(s)=\sum_{h=-r}^{r}\sum_{u=1}^{3}\frac{B_{h,u}}{(s-h)^u},
 \qquad
 B_{h,u}=\left.\frac1{(3-u)!}
 \left(\frac{d}{ds}\right)^{3-u}
 \bigl[(s-h)^3M_q(s)F(s)\bigr]\right|_{s=h}.
 \label{periods:k4-finite-residues}
\end{equation}
The complete rational coefficient is the finite expression
\begin{equation}
 R_q=-(r!)^4\sum_{h=-r}^{r}\sum_{u=1}^{3}B_{h,u}H_{r-h}^{(u)}.
 \label{periods:k4-rational-algorithm}
\end{equation}
All quantities in this formula are rational.  Its exact evaluation gives
\begin{equation}
 R_3=-\frac16,\qquad R_4=\frac7{360},\qquad
 R_5=\frac{1459}{60480},\qquad R_q>0\quad(6\le q\le20).
 \label{periods:k4-finite-base}
\end{equation}
The complete fractions for $3\le q\le20$, the source digests, and their
rational-arithmetic replay are part of the wheel certificate in
Appendix~\ref{app:certificates}.  Equations
\eqref{periods:principal-coefficients},
\eqref{periods:k4-finite-residues}, and
\eqref{periods:k4-rational-algorithm} specify every finite operation;
the finite check is used only for this bounded base range.

For a uniform proof beyond it, work on the six-variable unit simplex.
The polynomial $U_{K_4}$ contains sixteen of the twenty square-free
triple products.  Their complete sum is the third elementary symmetric
polynomial.  Averaging two variables with their sum fixed increases
this sum, since its only varying term is their product times the sum
of the other four variables.  Its maximum is therefore attained at
six equal variables, and
\begin{equation}
 U_{K_4}(t)\le\frac5{54},\qquad
 I_q(K_4;q-1,1,1;1,q-1,q-1)
 \ge\left(\frac{54}{5}\right)^q
 \frac{((q-2)!)^3}{(3q-1)!}.
 \label{periods:k4-integral-lower}
\end{equation}
After division by the positive zeta coefficient in
\eqref{periods:k4-value}, the lower bound is
\begin{equation}
 \left(\frac{54}{5}\right)^q
 \frac{(q-2)!(2q-2)!}{12(3q-1)!}.
 \label{periods:lower-ratio}
\end{equation}
At $q=21$ it equals
\begin{equation}
 \frac{2778146948724569113512799974260736}
 {2099879042434213173389434814453125}>\frac54,
 \label{periods:q21-ratio}
\end{equation}
because four times the numerator minus five times the denominator is
\[
 613192582727210587104025824777319>0.
\]
The ratio of
\eqref{periods:lower-ratio} at $q+1$ to its value at $q$ is
\begin{equation}
 \frac{36(q-1)(2q-1)}{5(3q+1)(3q+2)}>1\quad(q\ge6),
 \label{periods:ratio-growth}
\end{equation}
as the numerator minus denominator is
$27(q-6)^2+171(q-6)+80$.  Finally,
\begin{equation}
 \zeta(3)<1+\frac18+\int_2^{\infty}x^{-3}\,dx=\frac54.
 \label{periods:zeta3-bound}
\end{equation}
Equations \eqref{periods:k4-value}--\eqref{periods:zeta3-bound} imply
$R_q>0$ for every $q\ge21$.  Together with the exact finite base range,
this proves the lemma at every integer order.
\end{proof}

\subsection{The wheel space at a uniform sufficient threshold}
\label{periods:wheel-threshold}

\begin{theorem}[Complete wheel scalar span]
\label{periods:wheel-span}
For every $n\ge3$ and every integer $q\ge\max(3,n-1)$,
\begin{equation}
 \boxed{
 \mathcal P_q(W_n)=
 \mathbb Q+\sum_{j=2}^{n-1}\mathbb Q\,\zeta(2j-1).}
 \label{periods:wheel-space}
\end{equation}
The threshold is sufficient; optimality and arithmetic independence are
not asserted.
\end{theorem}

\begin{proof}
Let first $n\ge4$ and $q=n-1$.  Give the spokes the powers
\begin{equation}
 (q-1,q-1,\underbrace{q-2,\ldots,q-2}_{n-3},q-1)
 \label{periods:embedded-spokes}
\end{equation}
and the cyclic rim the powers $(1,\ldots,1,q-1)$.  Their sum is $nq$.
Eliminate rim vertices $3,4,\ldots,n-1$ in that order by
\eqref{periods:star-map}, combining every resulting parallel pair by
\eqref{periods:parallel-beta}.  At step $k$, vertex $k+2$ has incident
powers $(q-k-1,k,1)$, whose sum is $q$.  The two nontrivial parallel
factors are
\begin{equation}
 \frac1k,\qquad
 \begin{cases}
 1/[k(k+1)],&1\le k\le q-3,\\
 1/(q-2),&k=q-2.
 \end{cases}
 \label{periods:embedding-factors}
\end{equation}
Indeed, the first merge combines powers $q-k$ and $q-1$ into
$q-k-1$.  The second combines $q-2$ and $q-k$ at an intermediate
vertex, or $q-1$ and $q-k$ at the final vertex.  Formula
\eqref{periods:parallel-beta} gives exactly
\eqref{periods:embedding-factors}.  Their complete product is
$1/((q-2)!)^3$.

The remaining vertices are $0,1,2,n$, with powers
\begin{equation}
 \alpha_{01}=q-1,\qquad
 \alpha_{02}=\alpha_{0n}=\alpha_{12}=1,\qquad
 \alpha_{2n}=\alpha_{1n}=q-1.
 \label{periods:final-k4}
\end{equation}
This is the $K_4$ of Lemma~\ref{periods:k4-rational}.  All beta fibers
and the final graph converge, so the original positive integral
converges as well.  Its value is
\begin{equation}
 \frac{R_q}{((q-2)!)^3}
 +\frac{12}{(q-2)!(2q-2)!}\zeta(3),
 \qquad R_q\ne0.
 \label{periods:embedded-k4-value}
\end{equation}

For every $2\le q'<n$, multiply the numerator of the radial form
\eqref{periods:radial-master} by $U_{W_n}^{n-1-q'}$ and raise its
denominator to the same order.  Every resulting monomial converges by
\eqref{periods:dominated-shift}.  The triangular coefficients
\eqref{periods:radial-diagonal} therefore put every displayed odd zeta
value into $\mathcal P_{n-1}(W_n)$.  Subtract the zeta term from
\eqref{periods:embedded-k4-value} and divide by its nonzero rational
coefficient to obtain $1$ at the same level.

For $n=3$, use $q=3$ and spoke powers $(2,1,2)$ with rim powers
$(1,1,2)$.  The middle rim vertex is a unique star with powers $(1,1,1)$.
The star--triangle transformation and its parallel factors reduce it
to the one-loop triangle with powers $(1,1,1)$, whose period is $1/2$
by \eqref{periods:cycle-moments}.  The radial value at $q'=2$ is
$6\zeta(3)$ and lifts to $q=3$ by one multiplication by $U$.
This supplies both directions for $W_3$.

The upper inclusion is Proposition~\ref{periods:wheel-upper-proposition}.
Nesting \eqref{periods:dimension-shift} extends equality from the
indicated starting level to every larger integer $q$.
\end{proof}

As an additional normalized control, at $q=n$ give the first and last
spokes power $q-1$, all other spokes power $q-2$, the open rim path
power one, and the closing rim edge power $q-1$.  Eliminating vertices
$2,\ldots,n-1$ gives the same factor $1/((q-2)!)^3$, followed by a
one-loop triangle with powers $(1,1,q-2)$.  Its exact wheel period is
\begin{equation}
 \frac{1}{(n-1)(n-2)((n-2)!)^3}.
 \label{periods:rational-wheel-control}
\end{equation}
The accompanying wheel certificate checks this independent graph
reduction, the embedded $K_4$ construction, and complete radial coefficient
matrices through twelve loops.  These finite controls audit the formulas;
the all-order statements follow from the preceding proofs and the
explicit bounded base verification in
\eqref{periods:k4-finite-base}.

For a wheel core obtained after the complete network and series-chain
reduction, the rational factorial coefficients preserve
\eqref{periods:wheel-upper-space} coefficientwise in all retained source
spinors.  A specified source support or insertion family can span less
than the full scalar space.  Theorem~\ref{periods:wheel-span} does not
claim an all-order three-zeta closure for the triangular prism,
$K_{3,3}$, or $K_5$, nor a wheel threshold $q=3$ independent of $n$.
